\section{A stronger concentration result through local Lipschitz estimates}
\label{app:second_approach}

This appendix proves a concentration estimate for the normalized induced model.
We write
\[
    S_{HS}(d^2,s)=\{X\in M_{d^2,s}(\C):\norm{X}_2=1\}
\]
for the Hilbert--Schmidt unit sphere, viewed as a real Euclidean sphere of
dimension \(2d^2s\).  For
\[
    f_k(X)=\Tr\left(\left((XX^*)^\Gamma\right)^k\right),
    \qquad X\in S_{HS}(d^2,s),
\]
the key point is that \(f_k\) is not globally Lipschitz at the optimal scale,
but it is Lipschitz at that scale on a subset of overwhelming measure.  A
localized form of Levy's lemma then exploits the full sphere dimension,
which is of order \(d^4\) when \(s\sim\lambda d^2\).
The only non-elementary random-matrix input used below is Aubrun's
trace-moment estimate for the off-diagonal part of a partially transposed
Wishart matrix, stated explicitly as Theorem~\ref{lem:aubrun_moment_estimate}.

\begin{lemma}[Localized Levy lemma]
\label{lem:localized_levy}
Let \(S^{n-1}\) be the Euclidean unit sphere with its normalized Haar
measure.  Let \(\Omega\subset S^{n-1}\) be measurable and assume
\(\Prob(\Omega)\ge 3/4\).  Let \(h:S^{n-1}\to\R\) be measurable and assume
that \(h_{\mid\Omega}\) is \(L\)-Lipschitz.  If \(M_h\) is a median of \(h\),
then, for every \(t>0\),
\[
    \Prob\{|h-M_h|\ge t\}
    \le
    2\Prob(\Omega^c)+4\exp\left(-\frac{nt^2}{16L^2}\right).
\]
\end{lemma}

\begin{proof}
Define the Lipschitz extension
\[
    H(x)=\inf_{y\in\Omega}\bigl(h(y)+L\norm{x-y}_2\bigr).
\]
Then \(H\) is \(L\)-Lipschitz on the whole sphere and \(H=h\) on \(\Omega\).
Let \(M_H\) be a median of \(H\).  The spherical isoperimetric inequality says
that if \(A\) has spherical measure at least \(1/2\), then
\[
    \Prob\{\dist(x,A)>r\}\le \exp(-nr^2/4).
\]
Apply this to the median sublevel and superlevel sets of \(H\).  Since
\(|H(x)-H(y)|\le L\norm{x-y}_2\), the events
\(\{H>M_H+t\}\) and \(\{H<M_H-t\}\) are contained in the complements of the
\(t/L\)-neighborhoods of those two median sets.  The union of the two
one-sided bounds gives
\[
    \Prob\{|H-M_H|\ge t\}
    \le 2\exp\left(-\frac{nt^2}{4L^2}\right),
\]
with the displayed normalization of constants.  Since \(H=h\) on \(\Omega\),
\[
\begin{aligned}
    \Prob\{|h-M_H|\ge t\}
    &\le \Prob(\Omega^c)+\Prob\{|H-M_H|\ge t\}  \\
    &\le \Prob(\Omega^c)+2\exp\left(-\frac{nt^2}{4L^2}\right).
\end{aligned}
\]
This proves the estimate with the central value \(m_h=M_H\), with a slightly
stronger bad-set coefficient than stated for this auxiliary center.

We now replace \(M_H\) by a median \(M_h\) of \(h\).  We use the elementary
fact that, for any real number \(a\), any median \(M\) of a real random
variable \(Z\), and every \(t>0\),
\[
    \Prob\{|Z-M|\ge t\}
    \le 2\Prob\{|Z-a|\ge t/2\}.
\]
Indeed, if \(|M-a|<t/2\), then
\[
    \{|Z-M|\ge t\}\subset \{|Z-a|\ge t/2\}.
\]
If \(|M-a|\ge t/2\), then the median property implies
\(\Prob\{|Z-a|\ge t/2\}\ge1/2\), and the displayed inequality is trivial.
Applying this with \(Z=h\), \(M=M_h\), and \(a=M_H\), then using the estimate
already proved at scale \(t/2\), gives
\[
\begin{aligned}
    \Prob\{|h-M_h|\ge t\}
    &\le 2\Prob\{|h-M_H|\ge t/2\} \\
    &\le 2\Prob(\Omega^c)
       +4\exp\left(-\frac{nt^2}{16L^2}\right).
\end{aligned}
\]
\end{proof}

\begin{theorem}[Aubrun's moment estimate for the off-diagonal part]
\label{lem:aubrun_moment_estimate}
Let \(G\in M_{d^2,s}(\C)\) have independent standard complex Gaussian
entries, set
\[
    W=\frac1s GG^*,\qquad Y=W^\Gamma,\qquad Z=Y-\diag(Y).
\]
There exists a polynomial \(Q\), independent of \(d\) and \(s\), such that,
for every integer \(m\ge1\),
\[
    \left|\E\Tr(Z^m)\right|
    \le
    \left(\frac2s\right)^m
    (d+Q(m))^{m+2}(\sqrt{s}+Q(m))^m.
\]
In the proof below we only use this estimate for even \(m\), where \(Z^m\) is
positive semidefinite because \(Z\) is Hermitian.
\end{theorem}

\begin{proof}
Write a cyclic word in \(\{1,\ldots,d\}^2\) as
\[
    w=(i_0,i_1,\ldots,i_m),\qquad i_m=i_0,
\]
and write \(i_r=(a_r,b_r)\).  Expanding the trace gives
\[
    \Tr(Z^m)=\sum_w\prod_{r=0}^{m-1}Z_{i_r,i_{r+1}}.
\]
For an off-diagonal edge \(i_r\ne i_{r+1}\), the partial transpose gives
\[
    Z_{i_r,i_{r+1}}
    =
    \frac1s\sum_{\alpha_r=1}^s
    G_{(a_r,b_{r+1}),\alpha_r}\,
    \overline{G_{(a_{r+1},b_r),\alpha_r}}.
\]
For a diagonal edge, the corresponding entry of \(Z=Y-\diag(Y)\) is \(0\).
Hence
\[
\begin{aligned}
    \E\Tr(Z^m)
    &=
    s^{-m}\sum_w\mathbf 1_{\mathrm{off}}(w)
    \sum_{\alpha_0,\ldots,\alpha_{m-1}=1}^s
    \E\prod_{r=0}^{m-1}
    G_{(a_r,b_{r+1}),\alpha_r}\,
    \overline{G_{(a_{r+1},b_r),\alpha_r}} .
\end{aligned}
\]
The complex Wick formula turns the expectation into a count of permutations
\(\pi\in\mathfrak S_m\) for which every holomorphic coordinate is matched
with an equal anti-holomorphic coordinate:
\[
    (a_r,b_{r+1},\alpha_r)
    =
    (a_{\pi(r)+1},b_{\pi(r)},\alpha_{\pi(r)})
    \qquad(0\le r<m),
\]
with cyclic indices.  Let \(\mathcal F_\pi\) be the set of closed walks and
sample words satisfying these equalities.  Then
\[
    \E\Tr(Z^m)
    =
    s^{-m}\sum_{\pi\in\mathfrak S_m}
    \sum_{(w,\alpha)\in\mathcal F_\pi}\mathbf 1_{\mathrm{off}}(w),
\]
and therefore
\[
    \left|\E\Tr(Z^m)\right|
    \le
    s^{-m}\sum_{\pi\in\mathfrak S_m}\card\mathcal F_\pi.
\]

For a fixed surviving permutation, the matching equalities generate
equivalence relations among the left coordinates, right coordinates, and
sample labels.  If its profile is \((u,v)\), meaning \(u\) free left/right
classes and \(v\) free sample classes, then one may choose the free
left/right coordinates in at most \(d^u\) ways and the sample labels in at
most \(s^v\) ways.  Thus
\[
    \card\mathcal F_\pi\le d^u s^v.
\]
Regrouping by profiles gives
\[
    \sum_{\pi\in\mathfrak S_m}\card\mathcal F_\pi
    \le
    \sum_{u,v}N_m(u,v)d^u s^v.
\]
The Aubrun relation-counting estimate, proved by the innovation/defect
argument, states that the last weighted sum is bounded by
\[
    2^m(d+Q(m))^{m+2}(\sqrt{s}+Q(m))^m
\]
for a polynomial \(Q\) independent of \(d\) and \(s\).  Multiplying by
\(s^{-m}\) gives the announced estimate.
\end{proof}

\begin{lemma}[Standard Gaussian and gamma estimates]
\label{lem:standard_gaussian_gamma_estimates}
There is a universal constant \(C>0\) such that the following estimates hold.
\begin{enumerate}
    \item If \(G\in M_{m,n}(\C)\) has independent standard complex Gaussian
    entries, then
    \[
        \E\norm{G}_\infty\le C(\sqrt m+\sqrt n).
    \]
    \item If \(\xi_1,\ldots,\xi_N\) are random variables of the form
    \[
        \xi_i=\frac1s\sum_{\ell=1}^s |g_{i\ell}|^2,
    \]
    where the \(g_{i\ell}\)'s are independent standard complex Gaussian
    variables, then
    \[
        \E\max_{1\le i\le N}\xi_i
        \le
        1+C\sqrt{\frac{\log N}{s}}+C\frac{\log N}{s}.
    \]
\end{enumerate}
\end{lemma}

\begin{proof}
For the first estimate, fix \(1/4\)-nets \(N_m\) and \(N_n\) of the unit
spheres in \(\C^m\) and \(\C^n\), with \(|N_m|\le 9^{2m}\) and
\(|N_n|\le 9^{2n}\).  The net-lifting estimate gives
\[
    \norm G_\infty\le 2\sup_{u\in N_m,\ v\in N_n}|\langle u,Gv\rangle|.
\]
For fixed unit \(u,v\), the scalar \(\langle u,Gv\rangle\) is a standard
complex Gaussian, so
\[
    \Prob\{|\langle u,Gv\rangle|\ge r\}\le e^{-cr^2}
\]
for a universal \(c>0\).  The union bound over the two nets gives
\[
    \Prob\{\norm G_\infty\ge C(\sqrt m+\sqrt n)+t\}\le e^{-ct^2}.
\]
Integrating this tail in \(t\) gives
\[
    \E\norm G_\infty\le C'(\sqrt m+\sqrt n).
\]

For the second estimate, the gamma concentration bound
\[
    \Prob\left\{\xi_i\ge 1+2\sqrt{\frac u s}+2\frac u s\right\}\le e^{-u}
\]
is proved by the elementary Chernoff calculation for the MGF
\((1-\theta)^{-s}\) of a sum of \(s\) independent mean-one exponentials, with
\(\theta=\sqrt{u/s}/(1+\sqrt{u/s})\).  A union bound gives the same estimate for the
maximum with \(u\) replaced by \(u+\log N\).  Integrating this tail bound yields
the displayed expectation estimate after changing the universal constant.
\end{proof}

\begin{lemma}[Operator-norm inputs for the normalized model]
\label{lem:normalized_operator_norm_inputs}
Fix \(\lambda>0\), and let \(s=s_d\) satisfy \(s/d^2\to\lambda\).  Let \(X\)
be uniform on \(S_{HS}(d^2,s)\).  There exist constants \(C_1,C_2,c>0\),
depending only on \(\lambda\), such that, for all sufficiently large \(d\),
\[
    \E\norm{X}_\infty\le \frac{C_1}{d},
    \qquad
    \E\norm{(XX^*)^\Gamma}_\infty\le \frac{C_2}{d^2}.
\]
Furthermore, for suitable constants \(a,b>0\), depending only on \(\lambda\),
the sets
\[
    \Omega_1=\left\{X:\norm{X}_\infty\le \frac{a}{d}\right\},
    \qquad
    \Omega_2=\left\{X:\norm{(XX^*)^\Gamma}_\infty\le \frac{b}{d^2}\right\}
\]
satisfy
\[
    \Prob(\Omega_1^c)\le e^{-cd^2},
    \qquad
    \Prob(\Omega_2^c)\le e^{-cd^2}.
\]
\end{lemma}

\begin{proof}
Let \(G\in M_{d^2,s}(\C)\) have independent standard complex Gaussian entries.
Write
\[
    R=\norm{G}_2,\qquad X=\frac{G}{R}.
\]
Then \(X\) is uniform on \(S_{HS}(d^2,s)\) and is independent of \(R\).  Since
\(G=RX\), independence gives
\[
    \E\norm{G}_\infty=(\E R)(\E\norm{X}_\infty).
\]
The standard Gaussian operator-norm estimate
\[
    \E\norm{G}_\infty\le C(\sqrt{d^2}+\sqrt{s})
\]
from Lemma~\ref{lem:standard_gaussian_gamma_estimates}
and the elementary bound \(\E R\ge c\,d\sqrt{s}\) imply
\[
    \E\norm{X}_\infty
    \le
    \frac{C(d+\sqrt{s})}{c\,d\sqrt{s}}
    \le \frac{C_1}{d},
\]
because \(s\sim\lambda d^2\).

We next prove the corresponding expectation bound for the partial transpose.
Set
\[
    W=\frac1s GG^*,\qquad Y=W^\Gamma,
    \qquad Z=Y-\diag(Y).
\]
We apply Theorem~\ref{lem:aubrun_moment_estimate} only for even \(m\).  Since
\(Z\) is Hermitian,
\[
    \norm{Z}_\infty^m\le \Tr(Z^m)
\]
for even \(m\), and therefore
\[
\begin{aligned}
    \E\norm{Z}_\infty
    &\le \left(\E\norm{Z}_\infty^m\right)^{1/m}              \\
    &\le \left(\E\Tr(Z^m)\right)^{1/m}                       \\
    &\le
    \frac2s (d+Q(m))^{1+2/m}(\sqrt{s}+Q(m)).
\end{aligned}
\]
Choose an even integer \(m=m(d)\) such that
\[
    Q(m)=o(d),\qquad \log d=o(m).
\]
This is possible because \(Q\) is a fixed polynomial, for instance by taking
\(m=(\log d)^A\) with \(A\) large enough.  This is the same growth condition
used immediately after Aubrun's Proposition 7.1.  Since \(s\sim\lambda d^2\),
for large \(d\),
\[
\begin{aligned}
    \frac2s (d+Q(m))^{1+2/m}(\sqrt{s}+Q(m))
    &\le
    C_\lambda\, d^{2/m}                                      \\
    &=
    C_\lambda\, \exp\!\left(\frac{2\log d}{m}\right)
    \le 2C_\lambda.
\end{aligned}
\]
After changing the constant, this gives
\[
    \E\norm{Z}_\infty\le C_\lambda
\]
for all sufficiently large \(d\).

For the diagonal part, partial transposition preserves diagonal entries, so
\[
    Y_{ii}=W_{ii}=\frac1s\sum_{\ell=1}^{s}|G_{i\ell}|^2,
    \qquad 1\le i\le d^2.
\]
The variables \(Y_{ii}\) are of the form covered by
Lemma~\ref{lem:standard_gaussian_gamma_estimates}, with \(N=d^2\).  Hence
\[
    \E\max_{1\le i\le d^2}Y_{ii}
    \le
    1+C\sqrt{\frac{\log d}{s}}+C\frac{\log d}{s}
    \le C_\lambda.
\]
Consequently
\[
    \E\norm{Y}_\infty
    \le
    \E\norm{Z}_\infty+\E\norm{\diag(Y)}_\infty
    \le C_\lambda.
\]

Returning to the unnormalized Gaussian matrix, note that no concentration
estimate for the radial variable is needed at this point.  We have the exact
identity
\[
    (GG^*)^\Gamma=sY=R^2(XX^*)^\Gamma.
\]
Since \(R\) is independent of \(X\), and since \(\E R^2=d^2s\),
\[
    s\,\E\norm{Y}_\infty
    =
    \E\norm{(GG^*)^\Gamma}_\infty
    =
    (\E R^2)\,\E\norm{(XX^*)^\Gamma}_\infty
    =
    d^2s\,\E\norm{(XX^*)^\Gamma}_\infty.
\]
Thus
\[
    \E\norm{(XX^*)^\Gamma}_\infty
    =
    \frac{\E\norm{Y}_\infty}{d^2}
    \le \frac{C_2}{d^2}.
\]

We now pass from expectations to high-probability good sets.  The map
\[
    X\longmapsto \norm{X}_\infty
\]
is \(1\)-Lipschitz on the Hilbert--Schmidt sphere.  If \(M_1\) is a median of
\(\norm{X}_\infty\), then \(M_1\le 2\E\norm{X}_\infty\le 2C_1/d\).
Choosing \(a>2C_1+1\), Levy's lemma gives
\[
    \Prob(\Omega_1^c)\le e^{-cd^2}
\]
for large \(d\), since the sphere dimension is \(2d^2s\asymp d^4\) and the
deviation from the median is at least \(d^{-1}\).

It remains to control \(\Omega_2\).  Let
\[
    g(X)=\norm{(XX^*)^\Gamma}_\infty.
\]
For \(X,Y\in S_{HS}(d^2,s)\),
\[
\begin{aligned}
    |g(X)-g(Y)|
    &\le
    \norm{(XX^*)^\Gamma-(YY^*)^\Gamma}_\infty  \\
    &\le
    \norm{(XX^*)^\Gamma-(YY^*)^\Gamma}_2       \\
    &=
    \norm{XX^*-YY^*}_2                         \\
    &\le
    (\norm{X}_\infty+\norm{Y}_\infty)\norm{X-Y}_2.
\end{aligned}
\]
The last step follows from
\[
    XX^*-YY^*=(X-Y)X^*+Y(X-Y)^*
\]
and the Hilbert--Schmidt ideal inequalities.
Therefore \(g_{\mid\Omega_1}\) is \(2a/d\)-Lipschitz.  For large \(d\),
\(\Prob(\Omega_1)\ge3/4\).  Let \(M_g\) be a median of \(g\).  Since \(g\ge0\),
\[
    M_g\le 2\E g\le \frac{2C_2}{d^2}.
\]
Choose \(b>2C_2+1\).  Then
\[
    \{g>b/d^2\}
    \subset
    \{|g-M_g|>(b-2C_2)d^{-2}\}.
\]
Applying the localized Levy lemma to \(g\) on \(\Omega_1\), with
\[
    L=\frac{2a}{d},
    \qquad
    t=(b-2C_2)d^{-2},
    \qquad
    n=2d^2s\asymp d^4,
\]
gives
\[
    \Prob(\Omega_2^c)
    \le
    2\Prob(\Omega_1^c)
    +4\exp\left(-c d^2\right)
    \le e^{-c'd^2},
\]
after decreasing \(c'\) and changing constants.
\end{proof}

\begin{theorem}[Local Lipschitz concentration for partial transpose moments]
\label{thm:local_lipschitz_concentration}
Fix \(k\ge1\) and \(\lambda>0\), and let \((s_d)_{d\ge1}\) be an integer
sequence such that \(s_d/d^2\to\lambda\).  Let \(X\) be uniformly distributed
on
\[
    S_{HS}(d^2,s_d)=\{X\in M_{d^2,s_d}(\C):\norm{X}_2=1\}.
\]
Set
\[
    f_k(X)=\Tr\left(\left((XX^*)^\Gamma\right)^k\right).
\]
There exist constants \(c,C>0\), depending only on \(k\) and \(\lambda\), such
that, for all sufficiently large \(d\) and every \(\epsilon>0\),
\begin{equation}
\label{eq:local_lipschitz_concentration_theorem}
    \Prob\left(\abs{f_k(X)-\E f_k(X)}
    \geq \frac{\epsilon}{d^{2k-2}}\right)
    \leq C e^{-c d^2}
    +C\exp\left(-c\,\frac{d^4\epsilon^2}{k^2}\right).
\end{equation}
\end{theorem}

\begin{proof}
Write \(s=s_d\).  Since \(s/d^2\to\lambda\), all constants below are uniform
for sufficiently large \(d\), with dependence allowed on \(k\) and
\(\lambda\).  The real dimension of the ambient Hilbert--Schmidt sphere is
\[
    n=2d^2s\asymp d^4.
\]

\paragraph{1. Deterministic local Lipschitz estimate.}
Let
\[
    A=(XX^*)^\Gamma,\qquad B=(YY^*)^\Gamma.
\]
The telescoping identity
\[
    A^k-B^k=\sum_{j=0}^{k-1}A^j(A-B)B^{k-1-j}
\]
together with
\[
    |\Tr Z|\le\norm{Z}_1,
    \qquad
    \norm{UVW}_1\le\norm{U}_\infty\norm{V}_1\norm{W}_\infty
\]
gives
\[
    |\Tr(A^k)-\Tr(B^k)|
    \le
    k\max(\norm{A}_\infty,\norm{B}_\infty)^{k-1}\norm{A-B}_1.
\]
Since \(A-B=(XX^*-YY^*)^\Gamma\), partial transposition is an isometry for the
Hilbert--Schmidt norm, and \(\norm{M}_1\le d\norm{M}_2\) for matrices in
\(M_{d^2}(\C)\), we have
\[
    \norm{A-B}_1
    \le d\norm{A-B}_2
    =
    d\norm{XX^*-YY^*}_2.
\]
Moreover,
\[
    XX^*-YY^*=(X-Y)X^*+Y(X-Y)^*,
\]
and the Hilbert--Schmidt ideal inequality gives
\[
    \norm{XX^*-YY^*}_2
    \le
    (\norm{X}_\infty+\norm{Y}_\infty)\norm{X-Y}_2.
\]
Therefore, for all \(X,Y\in S_{HS}(d^2,s)\),
\begin{equation}
\label{eq:lipschitz_f_final_appendix_full_rigorous}
    |f_k(X)-f_k(Y)|
    \le
    dk\max\left(
      \norm{(XX^*)^\Gamma}_\infty,
      \norm{(YY^*)^\Gamma}_\infty
    \right)^{k-1}
    (\norm{X}_\infty+\norm{Y}_\infty)\norm{X-Y}_2.
\end{equation}

\paragraph{2. The good set.}
Let \(\Omega_1,\Omega_2\) be the sets from
Lemma~\ref{lem:normalized_operator_norm_inputs}, and set
\[
    \Omega=\Omega_1\cap\Omega_2.
\]
Then
\[
    \Prob(\Omega^c)
    \le \Prob(\Omega_1^c)+\Prob(\Omega_2^c)
    \le 2e^{-cd^2}.
\]
On \(\Omega\), estimate
\eqref{eq:lipschitz_f_final_appendix_full_rigorous} gives
\[
\begin{aligned}
    |f_k(X)-f_k(Y)|
    &\le
    dk\left(\frac{b}{d^2}\right)^{k-1}
    \frac{2a}{d}\norm{X-Y}_2  \\
    &\le
    \frac{C_k}{d^{2k-2}}\norm{X-Y}_2.
\end{aligned}
\]
Thus \(f_k\) is \(L_k\)-Lipschitz on \(\Omega\), with
\[
    L_k=\frac{C_k}{d^{2k-2}}.
\]
Equivalently, after enlarging \(C_k\), we may write
\[
    L_k=\frac{C(k,\lambda)\,k}{d^{2k-2}}.
\]

\paragraph{3. Median concentration.}
Apply Lemma~\ref{lem:localized_levy} to \(f_k\) on \(\Omega\).  For a median
\(M_k\) of \(f_k\), and after absorbing numerical constants into \(c\) and
\(C\), we get
\[
    \Prob\left\{
      \abs{f_k(X)-M_k}\ge \frac{u}{d^{2k-2}}
    \right\}
    \le
    4e^{-cd^2}
    +2\exp\left(
      -c\,\frac{d^4u^2}{k^2}
    \right).
\]
Indeed,
\[
    \frac{n(u/d^{2k-2})^2}{4L_k^2}
    \asymp
    \frac{d^4u^2}{k^2}.
\]

\paragraph{4. Replacing the median by the mean.}
We first record a deterministic bound on the range of \(f_k\).  Since partial
transposition preserves the Hilbert--Schmidt norm,
\[
    \norm{(XX^*)^\Gamma}_2=\norm{XX^*}_2\le \norm{XX^*}_1=\norm{X}_2^2=1.
\]
If \(k=1\), then
\[
    f_1(X)=\Tr((XX^*)^\Gamma)=\Tr(XX^*)=1.
\]
If \(k\ge2\), writing the eigenvalues of \((XX^*)^\Gamma\) as
\((\alpha_i)_i\), we have
\[
    |f_k(X)|
    =
    \left|\sum_i \alpha_i^k\right|
    \le
    \sum_i |\alpha_i|^k
    \le
    \left(\sum_i |\alpha_i|^2\right)^{k/2}
    \le 1.
\]
Thus \(|f_k|\le1\) for every \(k\ge1\).

Using the layer-cake identity and the preceding median concentration estimate,
\[
\begin{aligned}
    |\E f_k-M_k|
    &\le
    \E|f_k-M_k| \\
    &=
    \int_0^{2}\Prob\{|f_k-M_k|\ge t\}\,dt.
\end{aligned}
\]
Changing variables \(t=u/d^{2k-2}\) and using the displayed tail bound gives
\[
    |\E f_k-M_k|
    \le
    C e^{-cd^2}
    +C\frac{k}{d^{2k}},
\]
where the Gaussian integral contributes \(Ck/d^{2k}\), since the sphere
dimension is of order \(d^4\).

Let
\[
    \delta_\epsilon=\frac{\epsilon}{d^{2k-2}}.
\]
If
\[
    |\E f_k-M_k|\le \frac{\delta_\epsilon}{2},
\]
then
\[
    \{|f_k-\E f_k|\ge\delta_\epsilon\}
    \subset
    \{|f_k-M_k|\ge\delta_\epsilon/2\}.
\]
The median concentration estimate at half scale therefore gives
\[
    \Prob\{|f_k-\E f_k|\ge\delta_\epsilon\}
    \le
    C e^{-cd^2}
    +C\exp\left(-c\frac{d^4\epsilon^2}{k^2}\right).
\]

It remains only to cover the complementary case
\[
    |\E f_k-M_k|> \frac{\delta_\epsilon}{2}.
\]
By the estimate just proved for the mean--median gap, this can occur only if
\[
    \epsilon
    \le
    C\left(kd^{-2}+d^{2k-2}e^{-cd^2}\right).
\]
For all sufficiently large \(d\), this implies \(\epsilon\le C'k d^{-2}\).
Consequently
\[
    \exp\left(-c\frac{d^4\epsilon^2}{k^2}\right)
    \ge e^{-C''}.
\]
After enlarging the constant \(C\) in
\eqref{eq:local_lipschitz_concentration_theorem}, the right-hand side of
\eqref{eq:local_lipschitz_concentration_theorem} is then at least \(1\).
The trivial bound
\[
    \Prob\{|f_k-\E f_k|\ge\delta_\epsilon\}\le1
\]
closes this small-deviation branch.
\end{proof}

\begin{remark}[Why the normalized estimate is stronger]
The threshold proof in the main text only needs concentration of finitely many
moments at a scale large enough to identify the sign of the limiting Hankel
determinants.  The Gaussian polynomial concentration used there is well
suited for that purpose, but it gives a stretched-exponential tail of the form
\(\exp(-c d^{1/k})\) for relative deviations of the unnormalized Wishart
moments.  The estimate above is sharper in
the normalized induced model \(\rho=XX^*\), where
\[
    \Tr((\rho^\Gamma)^k)=f_k(X)
    \quad\text{is typically of order}\quad d^{-(2k-2)}.
\]
At this natural scale, the local Lipschitz argument gives
\[
    \Prob\left(
      |f_k(X)-\E f_k(X)|\ge \epsilon d^{-(2k-2)}
    \right)
    \le
    C e^{-c d^2}
    +C\exp\left(-c\,\frac{d^4\epsilon^2}{k^2}\right).
\]
Thus the typical part of the concentration reflects the full real dimension
of the Hilbert--Schmidt sphere, \(2d^2s\simeq 2\lambda d^4\); the only price
paid is the exceptional set, of probability \(e^{-cd^2}\), on which the
operator norms are not controlled.
\end{remark}
